\scanpage{50}
\margin{49}
- \mspage{74} - \hfill n$^{\mathrm{o}}$ 269

PROPOSITION 3 - \uline{Soit} $\Phi$ \uline{une forme} $\epsilon$-\uline{hermitienne vérifiant} (T) (resp. \uline{la forme bilinéaire associée à une forme quadratique non dégénérée} $Q$). \uline{Soient} $F_1$ \uline{et} $F_2$ \uline{deux sous-espaces totalement isotropes} (resp. \uline{totalement singuliers}) \uline{maximaux de} $E$, \uline{l'un des deux étant de dimension finie}. \uline{Posons} $F=F_1\cap F_2$. \uline{Soit} $S_i$ $(i=1,2)$ \uline{un supplémentaire de} $F$ \uline{dans} $F_i$ ; \uline{posons} $S=S_1+S_2$. \uline{Il existe alors deux sous espaces} $G$ \uline{et} $H$ \uline{de} $E$ \uline{tels que}

a) \uline{Les sous-espaces} $G+F$, $S$ \uline{et} $H$ \uline{sont non isotropes et deux à deux orthogonaux},

b) $E$ \uline{est somme directe de} $F,S,G$ \uline{et} $H$,

c) \uline{Il n'y a aucun vecteur isotrope} (resp. \uline{singulier}) \uline{non nul dans} $H$,

d) $G$ \uline{est totalement isotrope} (resp. \uline{totalement singulier}).

\uline{De plus}, $F_1$ \uline{et} $F_2$ \uline{sont tous deux de dimensions finies et on a}
\[
\begin{gathered}\dim F_1=\dim F_2,\quad\dim G=\dim F,\quad\dim S_1=\dim S_2,\quad\operatorname{codim}H={}\\=2\dim F_1.\end{gathered}
\]

Remarquons d'abord que, si $N$ est un sous espace totalement isotrope (resp. totalement singulier) maximal, alors tout vecteur isotrope (resp. singulier) $x$ orthogonal à $N$ est un élément de $N$, car sinon $N+Ax$ contredirait le caractère maximal de $N$.
Donc si, pour $i=1,2$, $x_i$ est un vecteur isotrope (resp. singulier) de $F_i^{\mathrm o}$, on a $x_i\in F_i$. D'autre part, si $y$ est un élément de $S_1$ orthogonal à $S_2$, il est orthogonal à $F_1$ puisque $F_1$ est totalement isotrope, donc à $F$, et par suite à $F_2=S_2+F$. Comme $y$ est isotrope (resp. singulier) et qu'il est orthogonal à $F_2$, on
